\documentclass[12pt,a4paper]{article}
\usepackage[utf8]{inputenc}
\usepackage{amsmath, amssymb, amsfonts}
\usepackage{graphicx}
\usepackage{fullpage}
\usepackage{enumitem}
\title{SOAL UAS IF1220 MATEMATIKA DISKRIT}
\author{Kelas 3E Kunugigaoka Junior High School}
\date{Waktu: 150 menit}

\begin{document}

\noindent \textbf{Nama: Kalyca Nathania Benedicta Manullang} \\
\noindent \textbf{NIM: 13524071}
\vspace{0.3cm}

\maketitle

\begin{center}
\large \textbf{BAGIAN A: SOAL MUDAH (EASY)}
\end{center}

\section*{Soal 1 Logika dan Penalaran}
Seorang detektif sedang menyelidiki kasus pencurian di sebuah toko buku. Terdapat empat tersangka: Andi ($A$), Budi ($B$), Cici ($C$), dan Dodi ($D$). Dari hasil investigasi diperoleh fakta-fakta berikut:
\begin{enumerate}
    \item Jika Andi bersalah, maka Budi juga bersalah.
    \item Jika Cici bersalah, maka Dodi juga bersalah.
    \item Andi bersalah atau Cici bersalah.
    \item Budi tidak bersalah.
\end{enumerate}
Dari fakta-fakta tersebut, tentukan apakah Dodi pasti bersalah? Gunakan hukum-hukum logika dan metode penarikan kesimpulan yang sahih untuk membuktikannya!

\section*{Soal 2 Himpunan}
Sebuah perusahaan teknologi mengadakan survei terhadap 120 karyawan mengenai penggunaan bahasa pemrograman. Data yang diperoleh sebagai berikut:
\begin{itemize}
    \item 65 karyawan menguasai Python
    \item 50 karyawan menguasai Java
    \item 45 karyawan menguasai C++
    \item 25 karyawan menguasai Python dan Java
    \item 20 karyawan menguasai Python dan C++
    \item 15 karyawan menguasai Java dan C++
    \item 10 karyawan menguasai ketiganya
\end{itemize}
\begin{enumerate}[label=(\alph*)]
    \item Gambarkan diagram Venn dari data tersebut!
    \item Berapa banyak karyawan yang tidak menguasai ketiga bahasa pemrograman tersebut?
\end{enumerate}

\section*{Soal 3 Relasi dan Fungsi}
Diberikan himpunan $A = \{2, 3, 4, 5\}$ dan relasi $R$ pada $A$ didefinisikan sebagai:
$$R = \{(2,2), (2,4), (3,3), (3,5), (4,4), (5,5)\}$$
\begin{enumerate}[label=(\alph*)]
    \item Apakah $R$ bersifat refleksif? Jelaskan!
    \item Apakah $R$ bersifat setangkup (simetris)? Jelaskan!
    \item Apakah $R$ bersifat menghantar (transitif)? Jelaskan!
\end{enumerate}

\section*{Soal 4 Teori Bilangan}
Dengan menggunakan Algoritma Euclidean, tentukan:
\begin{enumerate}[label=(\alph*)]
    \item $\gcd(180, 72)$
    \item $\operatorname{lcm}(180, 72)$
\end{enumerate}

\section*{Soal 5 Kombinatorika}
Sebuah restoran cepat saji menawarkan menu paket dengan pilihan:
\begin{itemize}
    \item 3 jenis minuman
    \item 5 jenis makanan utama
    \item 4 jenis makanan penutup
\end{itemize}
Berapa banyak cara memilih:
\begin{enumerate}[label=(\alph*)]
    \item Satu minuman, satu makanan utama, dan satu makanan penutup?
    \item Satu minuman \emph{atau} satu makanan utama (tidak keduanya)?
\end{enumerate}

\newpage
\begin{center}
\large \textbf{BAGIAN B: SOAL SEDANG (MEDIUM)}
\end{center}

\section*{Soal 6 Induksi Matematika}
Buktikan dengan induksi matematika bahwa untuk setiap bilangan bulat positif $n$,
$$\sum_{i=1}^{n} i \cdot 3^i = \frac{3}{4}\left((2n-1)3^n + 1\right)$$

\section*{Soal 7 Aljabar Boolean}
Sederhanakan ekspresi Boolean berikut menggunakan hukum-hukum aljabar Boolean:
$$F(x,y,z) = x'y'z + x'yz + xy'z + xyz$$
Tuliskan setiap langkah penyederhanaan beserta hukum yang digunakan!

\section*{Soal 8 Barisan, Sumasi, Rekursi, dan Relasi Rekurens}
Selesaikan relasi rekurens berikut:
$$a_n = 3a_{n-1} - 2a_{n-2}$$
dengan kondisi awal $a_0 = 1$ dan $a_1 = 4$.

\section*{Soal 9 Graf}
Perhatikan graf berikut dengan himpunan simpul $V = \{A, B, C, D, E\}$ dan himpunan sisi:
$$E = \{(A,B), (A,C), (B,C), (B,D), (C,D), (C,E), (D,E), (A,E)\}$$
\begin{enumerate}[label=(\alph*)]
    \item Gambarkan graf tersebut!
    \item Apakah graf tersebut memiliki lintasan Euler? Jelaskan!
    \item Apakah graf tersebut memiliki sirkuit Hamilton? Jelaskan!
\end{enumerate}

\section*{Soal 10 Kompleksitas Algoritma}
Diberikan potongan algoritma berikut:
\begin{verbatim}
for i := 1 to n do
    for j := 1 to i do
        for k := j to n do
            x := x + 1
        next k
    next j
next i
\end{verbatim}
\begin{enumerate}[label=(\alph*)]
    \item Tentukan $T(n)$ sebagai fungsi dari $n$ (jumlah operasi $x := x + 1$)!
    \item Nyatakan $T(n)$ dalam notasi Big-O!
\end{enumerate}

\newpage
\begin{center}
\large \textbf{BAGIAN C: SOAL SULIT (HARD)}
\end{center}

\section*{Soal 11 Pohon}
Sebuah pohon biner penuh (setiap simpul internal memiliki tepat 2 anak) memiliki 50 simpul daun.
\begin{enumerate}[label=(\alph*)]
    \item Berapa jumlah simpul internal pada pohon tersebut?
    \item Berapa jumlah total simpul pada pohon tersebut?
    \item Jika pohon tersebut memiliki tinggi 6 (dengan akar berada pada level 0), berapa jumlah maksimum dan minimum simpul internal yang mungkin? Jelaskan!
\end{enumerate}

\section*{Soal 12 Relasi dan Fungsi}
Diberikan fungsi $f: \mathbb{Z} \to \mathbb{Z}$ dengan $f(x) = 3x - 2$ dan $g: \mathbb{Z} \to \mathbb{Z}$ dengan $g(x) = x^2 + 1$.
\begin{enumerate}[label=(\alph*)]
    \item Tentukan $(f \circ g)(x)$ dan $(g \circ f)(x)$!
    \item Apakah $f$ merupakan fungsi injektif? Jelaskan!
    \item Apakah $f$ merupakan fungsi surjektif? Jelaskan!
    \item Apakah $g$ merupakan fungsi bijektif? Jelaskan!
\end{enumerate}

\section*{Soal 13 Teori Bilangan}
Gunakan Chinese Remainder Theorem (CRT) untuk menyelesaikan sistem kongruensi berikut:
$$x \equiv 4 \pmod{5}, \qquad x \equiv 2 \pmod{6}, \qquad x \equiv 3 \pmod{7}$$
Tentukan solusi $x$ terkecil yang memenuhi ketiga kongruensi tersebut!

\section*{Soal 14 Kombinatorika}
Sebuah kata sandi terdiri dari 10 karakter. Karakter dapat berupa huruf kapital (A-Z) atau angka (0-9). Setiap karakter boleh berulang. Namun, kata sandi harus memenuhi aturan berikut:
\begin{itemize}
    \item Minimal harus ada 3 huruf
    \item Minimal harus ada 3 angka
\end{itemize}
Berapa banyak kata sandi yang mungkin?

\section*{Soal 15 Anti-LLM $\boldsymbol{\bigstar}$}

Sebuah mesin penjual otomatis menerima koin pecahan \textbf{Rp2.000} dan \textbf{Rp5.000} saja. Mesin tidak memberikan kembalian. Harga minuman yang tersedia adalah Rp12.000, Rp14.000, dan Rp17.000.

Sistem logika mesin bekerja sebagai berikut:
\begin{itemize}
    \item Jika total uang yang dimasukkan sama dengan harga minuman, maka minuman keluar.
    \item Jika total uang melebihi harga, minuman keluar (tanpa kembalian).
    \item Jika total uang kurang dari harga, mesin terus menerima koin sampai cukup.
\end{itemize}

Terdapat 4 sensor pada mesin:
\begin{itemize}
    \item $P$: mendeteksi apakah ada koin Rp2.000 yang dimasukkan (bernilai 1 jika ada).
    \item $Q$: mendeteksi apakah ada koin Rp5.000 yang dimasukkan (bernilai 1 jika ada).
    \item $R$: mendeteksi apakah total uang $\ge$ Rp12.000 (bernilai 1 jika ya).
    \item $S$: mendeteksi apakah total uang $\ge$ Rp14.000 (bernilai 1 jika ya).
\end{itemize}

Sistem akan mengeluarkan minuman ($M = 1$) jika dan hanya jika:
\begin{center}
\textbf{total uang $\ge$ Rp12.000} \qquad \textbf{DAN} \qquad \textbf{jika total = Rp10.000, maka $M=0$.}
\end{center}
\begin{enumerate}[label=(\alph*)]
    \item Buatlah tabel kebenaran lengkap untuk semua kemungkinan kombinasi nilai kebenaran $P, Q, R, S$ dan tentukan nilai $M$ untuk setiap kombinasi!
    \item Tentukan ekspresi Boolean bentuk \textbf{Sum-of-Products (SOP) lengkap} untuk $M$!
    \item Sederhanakan ekspresi Boolean tersebut menggunakan \textbf{peta Karnaugh 4 variabel}! Tunjukkan pengelompokan yang Anda lakukan!
    \item Gambarkan rangkaian gerbang logika dari ekspresi yang telah disederhanakan!
\end{enumerate}

\end{document}